\section{Related Work}
\subsection{Dexterous Manipulation from Human Demonstrations} 
By harnessing human motion as the fuel of robot data, robots can acquire natural and versatile behaviors~\cite{argall2009survey,mandikal2022dexvip,shaw2023videodex,ye2023learning,qin2022dexmv,luo2024omnigrasp,ze2025twist,he2024learning}. In contrast to the typically scarce and costly teleoperation datasets, human motion data offers a more abundant and economically accessible resource for training robots. Recent advancements~\cite{zhou2025you,kim2025uniskill,lum2025crossing,dan2025x,wangmimicplay} have seen numerous studies leveraging human videos to extract features that assist in downstream policy learning. However, these efforts predominantly focus on transferring knowledge to robots with parallel-jaw end-effectors, often overlooking the necessity of bridging the embodiment gap and considering the complexity of high action space. In contrast, dexterous hands present greater challenges in modeling grasping postures and interactions with objects. Some recent works~\cite{kareer2024egomimic,hoque2025egodex,qiu2025humanoid} have employed customized, high-precision equipment to collect egocentric human videos for training robot policies, yet they neglect to model the complexity of hand-object interactions. Additionally, several recent studies~\cite{liu2025dextrack,chenobject,li2025maniptrans} utilize reinforcement learning to translate human hand motions into robotic behaviors, but these approaches are typically limited to simple, single-object tasks or lack closed-loop sim2real implementations. With a suite of innovative designs, \ourshort overcomes these shortcomings and equips the robot to perform various challenging, high–degree-of-freedom bimanual dexterous manipulation tasks. Moreover, we effectively transform diverse types of human motion into deployable robot policies.


\subsection{Vision-based Sim2real Manipulation}
Sim2real has achieved notable advancements in locomotion, with policies trained in simulation successfully deployed on quadruped or humanoid robots to perform agile motions~\cite{zhuang2023robot,zhuanghumanoid,rudin2025parkour,allshire2025visual,lei2024unio,qi2023hand}. In terms of manipulation, recent research has increasingly focused on leveraging simulation to generate or augment training data to obtain deployable visuomotor policies~\cite{wang2024cyberdemo,maddukuri2025sim,yuanlearning,huagensim2,wanggensim,xu2025dexsingrasplearningunifiedpolicy,zhang2025robustdexgrasp}. 
However, manipulation tasks, particularly those involving bimanual dexterous manipulation, present greater challenges for sim2real transfer due to the presence of high-frequency, fine-grained visual information.
Several approaches~\cite{lin2024learning,lin2024twisting,chen2023sequential,lin2025sim} have utilized depth cameras to extract object poses, shapes, and other relevant information, combining these with proprioceptive state data to facilitate sim2real transfer. While these methods can accurately obtain robot and object state information, the resulting policies often lack the ability to visually perceive environmental and object changes, which in turn forces manual specification of numerous state variables. Another line of works~\cite{yuanlearning,bjorck2025gr00t,singh2024dextrah} employ extensive domain randomization to train generalizable visual policies. However, achieving diversity in such randomization schemes typically involves labor-intensive engineering efforts. Recent studies~\cite{qureshi2024splatsim,li2024robogsim} have explored the use of Gaussian splatting~\cite{kerbl20233d} to achieve photorealistic rendering, but this approach requires manual scanning of entire scenes and training tailored to specific robot embodiments and environments, which limits scalability. 

Depth images inherently preserve object shape and spatial structure, enabling efficient sim2real transfer without the need to overcome variations in textures~\cite{dalal2024local,lumdextrah}. For instance, DextrAH-G~\cite{lumdextrah} utilizes depth maps for sim2real transfer but necessitates manually designing finely tuned and specific noise. Here, we also leverage depth images to facilitate egocentric sim2real transfer in an end-to-end fashion. In contrast to DextrAH-G, our approach adopts a more generalizable depth image augmentation strategy, which eliminates the need for camera-specific noise modeling. By applying our depth image processing method, we not only achieve strong semantic alignment between simulated and real-world observations but also obtain robust and generalizable visuomotor policies capable of handling complex manipulation tasks.


\subsection{Mobile Manipulation}
Mobile manipulation tasks further compound the challenges of robotic manipulation~\cite{yenamandra2023homerobot,fu2024mobile,huang2023skill,Liu_2024,zhi2025closedloopopenvocabularymobilemanipulation,wu2023tidybot}. The robot should typically complete a navigation phase before proceeding with the manipulation task. Many recent studies~\cite{Liu_2024,fang2023anygrasp,langsam2023,zhi2025closedloopopenvocabularymobilemanipulation} have adopted a modular mobile manipulation framework and utilized pre-built maps for robot mobile manipulation. OK-Robot~\cite{Liu_2024} uses a 3D map of the room built with iPhone and leverages an open-vocabulary object detector to perform open-vocabulary object navigation with the map. It also adopts AnyGrasp~\cite{fang2023anygrasp} combined with LangSam~\cite{langsam2023} to perform open-vocabulary grasping in the real world. COME-robot ~\cite{zhi2025closedloopopenvocabularymobilemanipulation} utilizes a global object map built by the robot for global-level perception. It leverages GPT-4V~\cite{gpt4v} to perform target object perception and task planning. However, these map-based methods are fundamentally constrained to small indoor spaces, as they become computationally intractable for large-scale environments like outdoor areas. They also struggle in feature-deprived settings or mixed indoor-outdoor navigation scenarios, where reliable localization becomes particularly challenging.

Another line of research explores end-to-end training paradigms that unify mobile navigation and manipulation into a single framework~\cite{jiang2025behavior,xiong2024adaptive,meng2025aim,yang2024harmonic}. However, such approaches not only demand a larger number of demonstrations or task-specific training, but also fail to leverage prior knowledge from navigation systems, thus constraining them to short-horizon tasks. To equip visuomotor policies with in-the-wild navigation capabilities, we incorporate the image-goal navigation foundation model that relies exclusively on RGB inputs and enables generalizable navigation, while being seamlessly integrated with the trained manipulation policy.




